% theme: decomposition
\MajorQuestion{物質の分解}
\SetRowSpaceQanda

\Instruction{化学変化と分解に関する次の問いに答えなさい。}
\QQ{もとの物質とは異なる別の物質ができる変化を何というか。}{化学変化}
\QQ{1種類の物質が、2種類以上の別の物質に分かれる化学変化を\NamedBlank{decomposition}という。}{分解}
\QQ{加熱したときに起こる\RefBlank{decomposition}を何というか。}{熱分解}
\QQ{物質Aが物質Bと物質Cに分かれるような変化は、化合と分解のどちらか。}{分解}
\QQ{物質を分解して、もとの物質とは別の物質ができる変化は、物理変化と化学変化のどちらか。}{化学変化}

\Instruction{炭酸水素ナトリウムの熱分解に関する次の問いに答えなさい。}
\QQ{加熱すると、炭酸ナトリウム、二酸化炭素、水に分解する白色の固体は何か。}{炭酸水素ナトリウム}
\QQ{炭酸水素ナトリウムを熱分解すると残る白色の固体は何か。}{炭酸ナトリウム}
\QQ{石灰水を白くにごらせる気体は何か。}{二酸化炭素}
\QQ{水の発生を確かめるとき、青色から桃色に変わる紙は何か。}{塩化コバルト紙}
\QQ{炭酸ナトリウムの水溶液で、フェノールフタレイン溶液は何色になるか。}{濃い赤色}

\Instruction{酸化銀の熱分解と金属に関する次の問いに答えなさい。}
\QQ{酸化銀を加熱すると発生する気体は何か。}{酸素}
\QQ{酸化銀を加熱したあとに残る、白色の固体は何か。}{銀}
\QQ{酸化銀のように、2種類以上の物質が結びついてできている物質を何というか。}{化合物}
\QQ{銀や銅のように、1種類の物質からできていて、それ以上分解できない物質を何というか。}{単体}
\QQ{金属をみがくと出る特有の光沢を何というか。}{金属光沢}

\Instruction{水の電気分解に関する次の問いに答えなさい。}
\QQ{電流を流して物質を分解することを何というか。}{電気分解}
\QQ{水の電気分解で、陰極に発生する気体は何か。}{水素}
\QQ{水の電気分解で、陽極に発生する気体は何か。}{酸素}
\QQ{水の電気分解で、発生する水素と酸素の体積比を答えなさい。}{水素：酸素＝2：1}
\QQ{純粋な水に電流を流れやすくするために加える物質は何か。}{水酸化ナトリウム}

\Instruction{塩化銅水溶液の電気分解と物質の種類に関する次の問いに答えなさい。}
\QQ{塩化銅水溶液の電気分解で、陰極に付着する物質は何か。}{銅}
\QQ{塩化銅水溶液の電気分解で、陽極に発生する気体は何か。}{塩素}
\QQ{塩化銅を電気分解すると、銅と何に分かれるか。}{塩素}
\QQ{水や酸化銀のように、2種類以上の物質からできている物質を何というか。}{化合物}
\QQ{空気や食塩水のように、複数の物質が混じり合ったものを何というか。}{混合物}
